\section{Introduction}
\label{sec:introduction}

Quantum algorithms for scientific computing have left the era of isolated
prototypes.  Since the quantum linear-system algorithm of Harrow, Hassidim,
and Lloyd~\cite{harrow2009hhl}, a mature body of work has unified linear
algebra under block-encodings and the quantum singular-value transformation
(QSVT)~\cite{gilyen2019qsvt,low2017qsp,martyn2021grand}, turned
Hamiltonian simulation into a family of interchangeable black-box
kernels~\cite{berry2015taylor,childs2021trotter,low2019qubitization}, and
produced a rapidly growing zoo of quantum solvers for ordinary and partial
differential equations: linear-combination-of-Hamiltonian-simulations (LCHS)
methods~\cite{an2026lchs}, Schr\"odingerization~\cite{jin2024schrodingerization},
contour-based matrix-function decompositions, Carleman
linearization for nonlinear systems~\cite{leyton2008ode,berry2014ode,liu2021pnas,liu2023carleman},
quantum finite-element and spectral
methods~\cite{montanaro2016fem,childs2020spectral}, and applications to fluid
dynamics~\cite{hu2024schrocircuits}.

The programming support available for expressing these algorithms has not kept
pace.  Today, a researcher who wishes to implement, compare, or compose such
algorithms faces a mismatch between the \emph{mathematical} structure of the
algorithms and the \emph{computational} structure of the tools:

\begin{enumerate}
\item \textbf{Algorithms are stated relative to access models, not circuits.}
  QLSS papers assume ``given a block-encoding of $A$''; ODE solvers assume
  ``given a block-encoding of the generator $G$ and a state-preparation
  oracle for the initial vector''; input models are specified by sparse-access
  oracles or quantum RAM (QRAM) with precise XOR semantics
  \cite{giovannetti2008qram,dervovic2018qlss}.  The concrete realization of
  these oracles---gate-based QROM, a physical QRAM, alias sampling, or an
  application-specific reversible circuit---is a separate decision that the
  algorithm text deliberately defers.  Mainstream quantum SDKs force the
  opposite order: an oracle must be a concrete circuit object before the
  algorithm can be written around it, and switching realization later means
  rewriting the algorithm.
\item \textbf{Programs are deep, symbolic, and large.}  Scientific-computing
  programs assemble solvers from primitives through many layers (input model
  $\rightarrow$ block-encoding $\rightarrow$ transform $\rightarrow$ solver
  $\rightarrow$ readout), with iteration counts and system sizes that make
  eager flattening to gate lists meaningless: a repetition count of $2^{40}$
  is a legitimate parameter in this regime.  What is needed is an intermediate
  representation that \emph{preserves} modular calls and static repetition,
  so that programs remain small, serializable, reviewable, and composable,
  and so that backends---not frontends---decide how far to expand.
\item \textbf{Correctness obligations are quantitative and compositional.}
  Block-encoding scale factors $\alpha$, adjoint and controlled variants of
  every subroutine, alias-freedom of operands, spectral-norm promises,
  dissipativity assumptions, and discretization caveats must all be tracked
  through composition.  Languages that model a quantum program as a flat list
  of gates have nowhere to attach this information; languages that track types
  but not these domain obligations lose them silently.
\end{enumerate}

This paper describes \pqlang, a Python-embedded quantum programming language in the embedded-domain-specific-language tradition~\cite{hudak1996building,liu2013chiselq}
designed for this workload.  \pqlang{} (currently at version 0.8.0) takes the
position that a quantum scientific-computing language is a
\emph{composition and bookkeeping discipline}, not a gate constructor.  Its
design is organized around three mechanisms:

\begin{description}
\item[Pluggable oracles (\secref{sec:oracles}).]  Algorithms are written
  against \emph{abstract oracles} declared with one of nine named paradigms
  (unitary, block-encoding, XOR database, state-preparation isometry, sparse
  access, phase oracle, reversible function, algorithm stage).  Abstract
  oracles are carried as \emph{open modules}---modules with a signature but no
  body---in a register-level intermediate representation (RIR).  A separate
  linking pass, \texttt{bind}, attaches concrete implementations, checking
  paradigm, register-type, scale-factor, and capability compatibility, and
  supports partial binding so that one open program can be linked against
  gate networks, QRAM resources, or application-specific circuits in any
  combination.  QRAM is a first-class typed resource, and alternative
  data-loading realizations of the same oracle slot are interchangeable.
\item[Algorithm protocols (\secref{sec:oracles}, \secref{sec:ode}).]
  Algorithm libraries declare their input models as Python \texttt{Protocol}
  interfaces (structural types) and expose solver families as small
  \emph{generator protocols} that map problem objects to quantum programs.
  Because protocols live in the library, not in the language core, new
  algorithms---including entirely new solver families---can be added without
  touching the grammar, the IR, or the serializer, while contract checking
  (with structured, machine-readable issue reports) remains uniform.
\item[Scientific-computing primitives (\secref{sec:primitives}).]  A library
  of block-encoding algebra (products, sums, linear combinations, tensor and
  Kronecker structures), PREPARE--SELECT decompositions with gate, QRAM, and
  alias-sampling realizations, a qubitization/QSVT/QSP stack with a
  self-verifying quantum-signal-processing phase engine, reversible
  fixed-point arithmetic with status flags, and a compiler from pure Python
  numerical functions (via their abstract syntax trees) to reversible
  modules.  On top of these, \pqlang{} assembles \emph{six distinct ODE
  solver families} (LCHS, Schr\"odingerization, contour-based matrix
  decomposition, Carleman linearization for nonlinear ODEs, implicit-Euler
  history systems solved by quantum linear-system solvers, and truncated
  Taylor kernels) and several PDE paths, including heat evolution, nonlinear
  polynomial PDEs via a homotopy-analysis linearization (QHAM), Fokker--Planck
  dynamics, and a quantum finite-volume method (QFVM) for the compressible
  Euler equations built from reversible Roe-flux arithmetic
  (\secref{sec:ode}, \secref{sec:casestudies}).
\end{description}

Underneath, the register-level IR is the architectural center
(\secref{sec:language}): programs are directed acyclic graphs of
\emph{modules} with typed quantum registers and resources; \texttt{Call},
\texttt{Repeat}, \texttt{Control}, and \texttt{Adjoint} nodes remain symbolic
through serialization and export; adjoint and controlled usage are checked
against compositionally inferred capabilities before any backend is invoked;
and the IR is deliberately unitary---measurement, reset, and classical
feedback are host-side concerns, planned as readout actions rather than IR
instructions.  The same RIR program can be validated while open, bound and
exported to a modular textual backend (OriginIR-ext), lowered to a strict
gate set, executed by a dependency-free reference simulator, or executed on
native register-level backends---including, since this work, a native RIR
interpreter shipped with PySparQ that expands the module graph at
interpretation time.

\paragraph{Contributions.}  This paper makes the following contributions:

\begin{itemize}
\item \textbf{A requirements analysis} that derives, from the structure of
  quantum scientific computing itself, eight concrete language requirements
  (deferred oracle realization, realization independence, structure
  preservation, capability tracking, symbolic scale, promise provenance,
  readout separation, and library extensibility), and identifies how existing
  languages fail them (\secref{sec:requirements}, \secref{sec:related}).
\item \textbf{A language design} meeting these requirements: an embedded
  builder over typed register views; a module-preserving register-level IR
  with an open-module oracle system, a linking pass, capability checking, and
  a canonical JSON serialization with schema
  (\secref{sec:language}, \secref{sec:oracles}).
\item \textbf{A protocol-based algorithm architecture} in which input models,
  adapters, and solver generators are first-class Python objects with
  structural contracts and machine-readable reports, and in which the same
  problem object is interchangeable between solver families and QLSS kernels
  (\secref{sec:oracles}, \secref{sec:ode}).
\item \textbf{A primitive and solver inventory} for quantum scientific
  computing, with detailed feature- and coverage-comparison tables against
  existing languages and frameworks (\secref{sec:primitives},
  \secref{sec:ode}, Tables~\ref{tab:qpl-comparison}
  and~\ref{tab:ode-coverage}).
\item \textbf{A numerical verification campaign on real backends}
  (\secref{sec:implementation}): 475 experiments in thirteen groups cover
  every algorithm family, each checked against a correctness oracle that is
  independent of the implementation, most of them evaluating the entire
  input domain in a single superposed execution.  The campaign runs on
  infrastructure contributed with this paper---an RIR-native PySparQ
  interpreter and UnifiedQuantum's \texttt{Circuit.to\_matrix}---alongside
  the pre-existing reference executor and adapter, cross-checks up to four
  independent implementations, separates implementation from method error,
  and surfaced two genuine defects, which are recorded openly with
  reproductions.  What remains open is the parameter-scaling axis, the
  defect repairs, and resource estimation---tracked explicitly rather than
  overclaimed.
\item \textbf{A compositional resource estimator}
  (\secref{sec:implementation}): a strict Toffoli+Clifford+T+QRAM compile
  and a per-(module, control-count) ledger that multiplies through symbolic
  \texttt{Repeat} nodes without expansion, reporting qubit, Toffoli,
  exact-$T$, pending-rotation, and QRAM query counts---the latter, to our
  knowledge, not previously reported as a first-class language
  metric---with scaling validated against theoretical expectations and with
  two lowering gaps (and counting) exposed by the accounting itself.
\end{itemize}

Everything described in this paper is implemented in the open-source \pqlang{}
package; all code listings are drawn from its doctested tutorials, examples,
and test suite.

\paragraph{Outline.}  \secref{sec:requirements} develops the requirements
analysis.  \secref{sec:related} surveys related languages and frameworks.
\secref{sec:language} presents the language surface and the RIR.
\secref{sec:oracles} presents the oracle and protocol system.
\secref{sec:primitives} presents the scientific-computing primitives.
\secref{sec:ode} presents the differential-equation stack, and
\secref{sec:casestudies} three case studies.  \secref{sec:implementation}
describes implementation and verification status: the evidence ladder, the
numerical-verification campaign with its per-family outcomes and recorded
defects, the resource-estimation study at the strict gate-set level, and
the deliberately reserved parts (validation at scale and the physical layer
of resource estimation).  \secref{sec:discussion} discusses limitations and future work.
